\documentclass[11pt]{amsart}
\usepackage{amsmath,amssymb,amsthm}
\usepackage{graphicx}
\usepackage{tikz-cd}
\usepackage{xcolor}
\usepackage{flows}
\usepackage{loom}

\theoremstyle{plain}
\newtheorem{theorem}{Theorem}[section]
\newtheorem{lemma}[theorem]{Lemma}
\newtheorem{proposition}[theorem]{Proposition}
\newtheorem{corollary}[theorem]{Corollary}
\theoremstyle{definition}
\newtheorem{definition}[theorem]{Definition}
\newtheorem{convention}[theorem]{Convention}
\newtheorem{example}[theorem]{Example}
\theoremstyle{remark}
\newtheorem{remark}[theorem]{Remark}

% A theorem-like environment the preamble does not declare with \newtheorem, so the directive below is what
% makes it a taxon: this is the one taxon configuration the source contract offers (book 5.5.1).
\newenvironment{caveat}{\par\medskip\noindent\textbf{Caveat.}\begingroup\itshape}{\endgroup\par\medskip}
% !LOOM environment: caveat = Caveat, remark

\title{Balanced flows on finite quivers}
\author{The loom showcase}


\ProvidesPackage{flows}[2026/09/19 macros for the showcase quilt]
\DeclareMathOperator{\supp}{supp}
\DeclareMathOperator{\rank}{rank}
\newcommand{\ZZ}{\mathbb{Z}}
\newcommand{\RR}{\mathbb{R}}
\newcommand{\Zcyc}{\mathcal{Z}}
\newcommand{\bd}{\partial}
\newcommand{\quiv}{Q}

\NeedsTeXFormat{LaTeX2e}
\ProvidesPackage{loom}[2026/09/16 v0.1 loom macros]
% loom.sty is MIT-licensed so that it can travel with a paper; see the LICENSE beside it in the loom distribution.

% Dependencies the text does not express as \ref or a matchable \cite.
% Prints nothing.
\providecommand{\uses}[1]{}

% Marks a statement or proof as mathematically incomplete.
% Prints nothing; an author may \renewcommand it to print in draft builds.
\providecommand{\incomplete}[1]{}

% Inputs a file with every sectioning command inside it shifted one level down.
% Composes when nested files nest files.
\providecommand{\nest}[1]{\begingroup
  \let\chapter\section
  \let\section\subsection
  \let\subsection\subsubsection
  \let\subsubsection\paragraph
  \let\paragraph\subparagraph
  \input{#1}\endgroup}
